
\documentclass[../../../physics-standard_questions_all.tex]{subfiles}


\begin{document}

図のように，抵抗R（抵抗値$R$），コンデンサC（容量$C$），およびスイッチSからなる回路があり，
Rの電流（図の左向き）を$I$，Cの電荷（上側極板が正）を$Q$と表す．
はじめSは開いており，$Q=Q_0$であった．$t=0$にSを閉じる．

\begin{enumerate}

    \item $Q=kQ_0\,(0 \leqq k \leqq 1)$となった瞬間の$I$を，$k$を用いて表せ．
    
    \item $I$，$Q$を，それぞれ$t$の関数として表せ．

\end{enumerate}

\vspace{0.2\baselineskip}
{\centering
    \includegraphics{\subfix{fig/fig_em_cb_10/fig_em_cb_10_00.pdf}}
\par}
\vspace{0.2\baselineskip}

\end{document}



